\documentclass[12pt]{article}

\usepackage[margin=1in]{geometry}
\usepackage{amsmath,amssymb,booktabs,mathtools}
\usepackage{natbib}
\usepackage[colorlinks=true,linkcolor=blue,urlcolor=blue,citecolor=blue]{hyperref}

\newcommand{\R}{\mathcal R}
\newcommand{\E}{\mathcal E}
\newcommand{\Lbar}{\bar L}
\newcommand{\obs}{\mathrm{obs}}
\newcommand{\ff}{\mathrm{ff}}

\title{Physical Travel Time in the Directed Route Model}
\author{}
\date{July 16, 2026}

\begin{document}
\maketitle

\section{Locked notation}

Three cost objects enter the model:
\begin{align*}
m_{kl}
&:= \text{physical travel time on directed link $(k,l)$, in minutes},\\
t_{kl}
&:=m_{kl}^{1/\theta}
   \quad\text{(link-level iceberg cost)},\\
\tau_{ij}
&:=\left[
  \sum_{r\in\R_{ij}}
  \prod_{(k,l)\in r}t_{kl}^{-\theta}
  \right]^{-1/\theta}
   \quad\text{(all-route bilateral cost index)}.
\end{align*}
The bilateral index $\tau_{ij}$ is not physical origin--destination travel
time. Its level therefore cannot be converted into seconds by multiplying by
60.

\section{Physical link time from the speed data}

Raw link length is measured in kilometres and raw speed is measured in
kilometres per hour. Physical time is consequently
\begin{equation}
  m_a^q
  =60\frac{\ell_a}{v_a^q},
  \qquad q\in\{\obs,\ff\},
  \label{eq:segment_time}
\end{equation}
where 60 converts hours into minutes. For an adjacent directed grid pair
$(k,l)$, segment minutes are added along the selected raw-road path $p_{kl}$:
\begin{equation}
  m_{kl}^q
  =\sum_{a\in p_{kl}}m_a^q.
  \label{eq:grid_time}
\end{equation}

\section{Physical time and iceberg cost}

For route $r\in\R_{ij}$, physical time remains additive:
\begin{equation}
  M_{ij,r}
  :=\sum_{(k,l)\in r}m_{kl}.
  \label{eq:route_minutes}
\end{equation}
Iceberg costs compound across links:
\begin{equation}
  C_{ij,r}
  :=\prod_{(k,l)\in r}t_{kl}
  =\prod_{(k,l)\in r}m_{kl}^{1/\theta}.
  \label{eq:route_iceberg}
\end{equation}
The product in \eqref{eq:route_iceberg} follows the iceberg-cost convention.
The transformation $t_{kl}=m_{kl}^{1/\theta}$ maps each link's minutes into
that convention; it does not replace the physical sum in
\eqref{eq:route_minutes}.

Conditional on origin $i$ and destination $j$, the route probability is
\begin{equation}
  \rho_{ij,r}
  =\frac{C_{ij,r}^{-\theta}}
  {\sum_{s\in\R_{ij}}C_{ij,s}^{-\theta}}.
  \label{eq:route_probability}
\end{equation}
Under the maintained mapping,
\begin{equation}
  C_{ij,r}^{-\theta}
  =\prod_{(k,l)\in r}m_{kl}^{-1}.
  \label{eq:theta_cancel}
\end{equation}
Equation \eqref{eq:theta_cancel} locates the role of the exponent. Recovering
minutes from a link-level iceberg cost requires $m_{kl}=t_{kl}^{\theta}$.
No additional $\theta$ power enters the additive route-time sum.

\section{Expected physical minutes conditional on an OD pair}

Expected physical time for an $i\to j$ commuter is
\begin{equation}
  M_{ij}
  :=\sum_{r\in\R_{ij}}\rho_{ij,r}M_{ij,r}.
  \label{eq:expected_od_minutes}
\end{equation}
This object differs from $\tau_{ij}$ and from $\tau_{ij}^{\theta}$.

The matrix representation avoids enumerating routes. Let the implemented
directed-edge kernel be
\begin{equation}
  A_{kl}
  =\sum_{e:k\to l}\frac{1}{\zeta m_e},
  \qquad
  R=(I-A)^{-1},
  \label{eq:resolvent}
\end{equation}
where $\zeta$ is the common route-diffusion normalization used to satisfy the
spectral-radius condition. Define
\begin{equation}
  B_{kl}
  =\sum_{e:k\to l}A_e m_e.
  \label{eq:reward_matrix}
\end{equation}
The numerator of expected route minutes sums each link-minute reward over all
routes:
\begin{equation}
  M_{ij}
  =\frac{(RBR)_{ij}}{R_{ij}}.
  \label{eq:matrix_minutes}
\end{equation}

\section{Aggregation across commuters}

Let $L_{ij,r}$ denote commuters choosing residence $i$, workplace $j$, and
route $r$. Directed link traffic is the expected number of link passages:
\begin{equation}
  \Xi_{kl}
  =\sum_{i,j}\sum_{r\in\R_{ij}}
  L_{ij,r}\,N_{kl}(r),
  \label{eq:link_passages}
\end{equation}
where $N_{kl}(r)$ counts the number of times route $r$ traverses directed
link $(k,l)$. Total person-minutes can be aggregated by commuter-route or by
directed link:
\begin{align}
  \sum_{i,j}\sum_{r\in\R_{ij}}L_{ij,r}M_{ij,r}
  &=\sum_{k,l}\Xi_{kl}m_{kl}.
  \label{eq:person_minutes_identity}
\end{align}
Average physical commuting time is therefore
\begin{equation}
  \bar M
  =\frac{1}{\Lbar}\sum_{k,l}\Xi_{kl}m_{kl}.
  \label{eq:average_minutes}
\end{equation}
For scenario $h\in\{ST,LT\}$, seconds saved per commuter and total
person-hours saved during one AM peak are
\begin{align}
  S_h^{\mathrm{sec}}
  &=-60\left(\bar M^h-\bar M^0\right),
  \label{eq:seconds_saved}\\
  S_h^{\mathrm{ph}}
  &=-\frac{\Lbar}{60}\left(\bar M^h-\bar M^0\right).
  \label{eq:person_hours_saved}
\end{align}

\section{Short-run and long-run decomposition}

The implemented link-passage decomposition is exact. The short-run term is
the complete change in physical time after route traffic and congestion
re-equilibrate with residence and workplace marginals fixed:
\begin{equation}
  D^{ST}=\bar M^{ST}-\bar M^0.
  \label{eq:st_change}
\end{equation}
The subsequent long-run adjustment is separated into a link-time effect at
short-run passages and a passage-reallocation effect evaluated at long-run
minutes:
\begin{align}
  G^{LT}
  &=\frac{1}{\Lbar}\sum_{k,l}
    \Xi_{kl}^{ST}\left(m_{kl}^{LT}-m_{kl}^{ST}\right),
  \label{eq:lt_network}\\
  R^{LT}
  &=\frac{1}{\Lbar}\sum_{k,l}
    m_{kl}^{LT}\left(\Xi_{kl}^{LT}-\Xi_{kl}^{ST}\right).
  \label{eq:lt_reallocation}
\end{align}
These terms satisfy
\begin{equation}
  \bar M^{LT}-\bar M^0=D^{ST}+G^{LT}+R^{LT}.
  \label{eq:decomposition_identity}
\end{equation}

\section{Mainline results}

Table \ref{tab:physical_results} reports the top-100 tidal-lane
counterfactual under the anchored-congestion mainline.
\begin{table}[ht]
\centering
\caption{Expected physical commuting time}
\label{tab:physical_results}
\begin{tabular}{lrrr}
\toprule
 & Baseline & Short run & Long run \\
\midrule
Average physical time (minutes) & 11.298 & 10.596 & 11.538 \\
Change from baseline (minutes) & 0.000 & $-0.702$ & 0.240 \\
Seconds saved per commuter & 0.0 & 42.1 & $-14.4$ \\
Person-hours saved per AM peak & 0.0 & 1,031 & $-352$ \\
Expected-utility welfare change (percent) & 0.00 & 0.87 & 0.32 \\
\bottomrule
\end{tabular}
\end{table}

The short-run physical-time gain is 42.1 seconds per commuter. Additional
network feedback contributes 26.6 seconds and long-run link-passage
reallocation contributes 30.0 seconds in the opposite direction. Average
physical time is consequently 14.4 seconds above baseline in the long run,
although expected utility remains 0.32 percent above baseline. Welfare also
loads productivity, amenities, wages, and access, so its sign need not match
the sign of average physical travel time.

\section{Implementation and numerical check}

Headline physical-time results are computed from
\eqref{eq:average_minutes} using the model's equilibrium arrays
\texttt{Xi\_kl} and \texttt{tau\_min\_kl}. The variable
\texttt{tau\_min\_kl} is $m_{kl}$ despite its legacy code name.

The solver floors very small entries of $R=(I-A)^{-1}$ at a numerical
constant. Reconstructing every $M_{ij}$ from the floored matrix through
\eqref{eq:matrix_minutes} therefore differs slightly in counterfactual
changes and more substantially in levels for remote OD pairs. The reported
minutes exploit the link-passage identity in \eqref{eq:average_minutes}; the
OD-matrix reconstruction is retained as a diagnostic until the resolvent
floor is revised and the equality in \eqref{eq:person_minutes_identity} is
restored numerically.

\end{document}
